% GENERATED FILE. Edit canonical lessons, then run npm run generate:books.
% canonical-source-hash: fnv1a64:d21d0b48764248af
% canonical-lessons: ZH-C309-qiangjie, ZH-C309-xiaotou, ZH-W309-tou-steal, ZH-C309-pian, ZH-C309-pianzi, ZH-C309-guyi

\chapter{A Thief}
\label{ch:zh-a-thief}
\hlchaptermodality{309}

\begin{chapteropening}
\textbf{By the end of this chapter:} \emph{I can talk about a robbery, a thief, a con man and something done on purpose.}

Complete the last lesson of chapter 309: I can talk about a robbery, a thief, a con man and something done on purpose.
\end{chapteropening}

\section[qiǎngjié]{qiǎngjié — to rob by force}

\label{lesson:ZH-C309-qiangjie}



\begin{quote}
Before the new one: say the Mandarin for to commit a crime, then the Mandarin for the local police station.
\end{quote}

\subsection*{You'll want to know: qiǎngjié}
\emph{qiǎngjié} — \textquotedblleft{}to rob by force\textquotedblright{}.

\begin{quote}
\emph{yǒu rén qiǎngjié le zhè jiā yínháng} — Someone robbed this bank.
\end{quote}

Its characters come later; for now it is a word for the ear.

\begin{cousinweb}
\emph{qiǎngjié} is built from \emph{qiǎng} 'to snatch' and \emph{jié} 'to plunder', both inherited from Middle Chinese, \emph{tshjangX} and \emph{kjaep} (Guangyun, in Baxter's transcription), which glosses the second 'to rob by force'; the sense 'to snatch' for the first is not given there.
\end{cousinweb}

\begin{grammarlens}[title={What you've built}]
One more word for telling what happened.
\end{grammarlens}

\subsection*{Your turn}
\begin{itemize}
\raggedright
  \item \emph{Say it:} \emph{qiǎngjié}
  \item \emph{Say it:} \emph{qiǎngjié}, once more
  \item \emph{Say it:} \emph{yǒu rén qiǎngjié le zhè jiā yínháng} — Someone robbed this bank.
  \item \emph{Recall:} say the Mandarin for to commit a crime, then the Mandarin for the local police station, then say \emph{qiǎngjié} again
\end{itemize}

\begin{tcolorbox}[breakable,colback=teal!4,colframe=teal!35!black,title={Before you move on}]
What does qiǎngjié mean? (\textquotedblleft{}To rob by force\textquotedblright{}.) Say it once more.
\end{tcolorbox}
\section[\texorpdfstring{\zh{小偷} (xiǎotōu)}{小偷 (xiǎotōu)}]{\zh{小偷} (xiǎotōu) — thief}

\label{lesson:ZH-C309-xiaotou}



\begin{quote}
Before the new one: say the Mandarin for the local police station, then the Mandarin for to rob by force.
\end{quote}

\subsection*{You'll want to know: \zh{小偷}}
\textbf{\zh{小偷}} — \emph{xiǎotōu} — \textquotedblleft{}thief\textquotedblright{}.

\begin{quote}
\textbf{\zh{小偷偷了我的钱包}} \emph{xiǎotōu tōu le wǒ de qiánbāo} — A thief stole my wallet.
\end{quote}

The new character, in pieces: \zh{偷} — eleven strokes in two pieces: the person piece of \zh{你} on the left, two strokes; on the right a roof of two strokes with a short bar under it, then \zh{月} on the left and the two strokes on the right of \zh{到}. The next lesson puts it in your hand.

\begin{cousinweb}
\textbf{\zh{小偷}} (\emph{xiǎotōu}) is built from \zh{小} 'small' and \zh{偷} 'to steal', both inherited from Middle Chinese, \emph{sjewX} and \emph{thuw} (Guangyun, in Baxter's transcription), which glosses the second 'to rob'.
\end{cousinweb}

\begin{grammarlens}[title={What you've built}]
One more word for telling what happened.
\end{grammarlens}

\subsection*{Your turn}
\begin{itemize}
\raggedright
  \item \emph{Say it:} \emph{xiǎotōu}
  \item \emph{Say it:} \emph{xiǎotōu}, once more
  \item \emph{Say it:} \emph{xiǎotōu tōu le wǒ de qiánbāo} — A thief stole my wallet.
  \item \emph{Recall:} say the Mandarin for the local police station, then the Mandarin for to rob by force, then say \emph{xiǎotōu} again
\end{itemize}

\begin{tcolorbox}[breakable,colback=teal!4,colframe=teal!35!black,title={Before you move on}]
What does \zh{小偷} mean? (\textquotedblleft{}Thief\textquotedblright{}.) Say it once more.
\end{tcolorbox}
\section[\texorpdfstring{\zh{偷} (tōu)}{偷 (tōu)}]{\zh{偷} — the person piece, then a roof over \zh{月} and two posts}

\label{lesson:ZH-W309-tou-steal}



\begin{quote}
The word first, then the character inside it.

\begin{itemize}
\raggedright
  \item \emph{Recall:} say \emph{xiǎotōu} — \textbf{R1}, one lesson back, before it decays
\end{itemize}

\begin{quote}
\zh{你} \zh{到}
\end{quote}

\emph{Read it:} each one aloud, with nothing said first — \emph{nǐ}, \emph{dào}
\end{quote}

\subsection*{Script — the person piece, then a roof over \zh{月} and two posts}
\hlblockfigure{\detokenize{figures/ZH-W309-tou-steal-filmstrip.pdf}}{How \zh{偷} is written, stroke by stroke}

\begin{quote}
\zh{偷}
\end{quote}

\textbf{Two pieces, left before right.} Count the pieces, not the strokes.

\textbf{The left piece, two strokes.} The person piece of \zh{你}: a short sweep down to the left, then a vertical.

\textbf{The right piece, nine strokes, top to bottom.} A roof: a sweep down to the left and a long stroke falling to the right. A short bar under it. Then, side by side, \zh{月} on the left, four strokes, and on the right the two strokes on the right of \zh{到}: a short vertical and a long vertical that hooks to the left.

2 + 9 = eleven.

\subsection*{Your turn}
\begin{itemize}
\raggedright
  \item \emph{Write it:} \zh{偷}, all eleven, with the model in front of you
  \item \emph{Cover:} write it again from memory, and count the pieces, not the strokes
\end{itemize}

\begin{tcolorbox}[breakable,colback=teal!4,colframe=teal!35!black,title={Before you move on}]
How many pieces in \zh{偷}? (\textbf{Two}.) What is under the roof and the bar? (\textbf{\zh{月}, and the two strokes on the right of \zh{到}}.) Which stroke is last? (\textbf{The long vertical that hooks to the left}.)

Source: \href{https://raw.githubusercontent.com/chanind/hanzi-writer-data/68d10a4b21150cae5e1ebbd223eed289cf32d90c/data/\%E5\%81\%B7.json}{Hanzi Writer Data \zh{偷}, pinned stroke paths and medians}.
\end{tcolorbox}
\section[piàn]{piàn — to cheat; to con}

\label{lesson:ZH-C309-pian}



\begin{quote}
Before the new one: say the Mandarin for to rob by force, then the Mandarin for a thief.
\end{quote}

\subsection*{You'll want to know: piàn}
\emph{piàn} — \textquotedblleft{}to cheat; to con\textquotedblright{}.

\begin{quote}
\emph{tā piàn le wǒ de qián} — He conned me out of my money.
\end{quote}

Its characters come later; for now it is a word for the ear.

\begin{cousinweb}
\emph{piàn} No dictionary we use records an older form, so its history is not settled.
\end{cousinweb}

\begin{grammarlens}[title={What you've built}]
One more word for telling what happened.
\end{grammarlens}

\subsection*{Your turn}
\begin{itemize}
\raggedright
  \item \emph{Say it:} \emph{piàn}
  \item \emph{Say it:} \emph{piàn}, once more
  \item \emph{Say it:} \emph{tā piàn le wǒ de qián} — He conned me out of my money.
  \item \emph{Recall:} say the Mandarin for to rob by force, then the Mandarin for a thief, then say \emph{piàn} again
\end{itemize}

\begin{tcolorbox}[breakable,colback=teal!4,colframe=teal!35!black,title={Before you move on}]
What does piàn mean? (\textquotedblleft{}To cheat; to con\textquotedblright{}.) Say it once more.
\end{tcolorbox}
\section[piànzi]{piànzi — swindler; con man}

\label{lesson:ZH-C309-pianzi}



\begin{quote}
Before the new one: say the Mandarin for a thief, then the Mandarin for to con someone.
\end{quote}

\subsection*{You'll want to know: piànzi}
\emph{piànzi} — \textquotedblleft{}swindler; con man\textquotedblright{}.

\begin{quote}
\emph{nàge rén shì piànzi} — That man is a con man.
\end{quote}

Its characters come later; for now it is a word for the ear.

\begin{cousinweb}
\emph{piànzi} is built from \emph{piàn} 'to con' and the noun ending \emph{zi}. The second is inherited from Middle Chinese \emph{tsiX} (Guangyun, in Baxter's transcription); the Guangyun has no entry for the first, so its history is not settled.
\end{cousinweb}

\begin{grammarlens}[title={What you've built}]
One more word for telling what happened.
\end{grammarlens}

\subsection*{Your turn}
\begin{itemize}
\raggedright
  \item \emph{Say it:} \emph{piànzi}
  \item \emph{Say it:} \emph{piànzi}, once more
  \item \emph{Say it:} \emph{nàge rén shì piànzi} — That man is a con man.
  \item \emph{Recall:} say the Mandarin for a thief, then the Mandarin for to con someone, then say \emph{piànzi} again
\end{itemize}

\begin{tcolorbox}[breakable,colback=teal!4,colframe=teal!35!black,title={Before you move on}]
What does piànzi mean? (\textquotedblleft{}Swindler; con man\textquotedblright{}.) Say it once more.
\end{tcolorbox}
\section[\texorpdfstring{\zh{故意} (gùyì)}{故意 (gùyì)}]{\zh{故意} (gùyì) — on purpose; deliberately}

\label{lesson:ZH-C309-guyi}



\begin{quote}
Before the new one: say the Mandarin for to con someone, then the Mandarin for a con man.
\end{quote}

\subsection*{You'll want to know: \zh{故意}}
\textbf{\zh{故意}} — \emph{gùyì} — \textquotedblleft{}on purpose; deliberately\textquotedblright{}.

\begin{quote}
\textbf{\zh{他故意不说话}} \emph{tā gùyì bù shuōhuà} — He's keeping quiet on purpose.
\end{quote}

\begin{cousinweb}
\textbf{\zh{故意}} (\emph{gùyì}) is built from \zh{故} and \zh{意} 'an intention', both inherited from Middle Chinese, \emph{kuH} and \emph{'iH} (Guangyun, in Baxter's transcription), which glosses the first 'old' and 'a matter' and the second 'will'; how the pair came to mean 'on purpose' is not settled.
\end{cousinweb}

\begin{grammarlens}[title={What you've built}]
One more word for telling what happened.
\end{grammarlens}

\subsection*{Your turn}
\begin{itemize}
\raggedright
  \item \emph{Say it:} \emph{gùyì}
  \item \emph{Say it:} \emph{gùyì}, once more
  \item \emph{Say it:} \emph{tā gùyì bù shuōhuà} — He's keeping quiet on purpose.
  \item \emph{Recall:} say the Mandarin for to con someone, then the Mandarin for a con man, then say \emph{gùyì} again
\end{itemize}

\begin{tcolorbox}[breakable,colback=teal!4,colframe=teal!35!black,title={Before you move on}]
What does \zh{故意} mean? (\textquotedblleft{}On purpose; deliberately\textquotedblright{}.) Say it once more.
\end{tcolorbox}
